\documentclass[a4paper,10pt]{article}
\usepackage[margin=1.4in]{geometry}
\usepackage[swedish]{babel}
\usepackage[utf8]{inputenc}
\usepackage{titlesec}
\usepackage{titling}
\usepackage{todonotes}



\setlength{\parskip}{1em}
\setlength{\parindent}{0pt}
\titlespacing{\section}{0pt}{\parskip}{-\parskip}
\titlespacing{\subsection}{0pt}{\parskip}{-\parskip}
\titlespacing{\subsubsection}{0pt}{\parskip}{-\parskip}
\titlespacing{\part}{0pt}{\parskip}{-\parskip}

\def\ftitle{Testrapport}
\def\fversion{0.1}

\begin{document}
\begin{titlepage} % Suppresses displaying the page number on the title page and the subsequent page counts as page 1
	\newcommand{\HRule}{\rule{\linewidth}{0.5mm}} % Defines a new command for horizontal lines, change thickness here

	\center % Centre everything on the page

	%------------------------------------------------
	%	Headings
	%------------------------------------------------

	\textsc{\LARGE Linköpings universitet \\ \vspace{0.2em} Institutionen för datavetenskap }\\[2cm]

    \large\today

    \vspace{1cm}


	%------------------------------------------------
	%	Title
	%------------------------------------------------

	\HRule\\[0.4cm]

	{\huge\bfseries Schemaläggningsstöd för  kirurgi \vspace{.1em} \\ \ftitle }\\[0.4cm] % Title of your document

	\HRule\\[1cm]

	%------------------------------------------------
	%	Author(s)
	%------------------------------------------------

	\begin{minipage}{0.7\textwidth}
			\large
            \emph{Version: \fversion}
            \vspace{1em}

            \textbf{\\Adam Andersson, Niclas Byrsten, \\Björn Hvass, Henrik Lindström, \\Martin Persson, Christoffer Sjöbergsson, \\Tor Utterborn}


            \vspace{1em}

            Handledare: Jonas Wallgren

            Examinator: Kristian Sandahl
	\end{minipage}
	~

	%------------------------------------------------
	%	Logo
	%------------------------------------------------

	%\vfill\vfill
	%\includegraphics[width=0.2\textwidth]{../Templates/Aeon}\\[1cm] % Include a department/university logo - this will require the graphicx package

	%----------------------------------------------------------------------------------------

	\vfill % Push the date up 1/4 of the remaining page

\end{titlepage}


\section*{\begin{center}Projektidentitet\end{center}}
    \vspace{-2.5em}
    \begin{center}
        \begin{tabular}{|c c c|}
        \hline
        Namn & Roll & E-post \\
        \hline
        Adam Andersson& Teamleader & adam.e.a.andersson@gmail.com\\
        \hline
        Niclas Byrsten & Testansvarig & nicby889@student.liu.se\\
        \hline
        Björn Hvass & Konfigurationsansvarig & bjorn.hvass@gmail.com\\
        \hline
        Henrik Lindström & Utvecklingsledare & henli070@student.liu.se\\
        \hline
        Martin Persson & Arkitekturansvarig & marpe902@liu.student.se\\
        \hline
        Christoffer Sjöbersson & Analysansvarig & chrsj812@liu.student.se\\
        \hline
        Tor Utterborn & Dokument- \& Kvalitetsansvarig & tor.utterborn@gmail.com\\
        \hline
        \end{tabular}
    \end{center}

    \begin{center}
        \small
        \textbf{Kund}\\Region Östergötland, 581 91 Linköping.

        \textbf{Kontaktperson hos kund}\\
        Gunnar Nordqvist, IT-arkitekt, 010-1030698, Gunnar.Nordqvist@regionostergotland.se\\
        Erik Sundvall, Informationsarkitekt, 010-1036252, Erik.Sundvall@regionostergotland.se
    \end{center}

\vspace{9em}



\section*{\begin{center}Dokumentationshistorik\end{center}}
    \begin{center}
        \begin{tabular}{|c c c c |}
        \hline
        Datum & händelse & iteration & version\\
        \hline
        2018-02-27 & Dokument skapas & 2 &  0.1\\
        \hline
        \end{tabular}
    \end{center}
    \clearpage
    \tableofcontents
    \clearpage

\section{Introduktion}
Nedan så listas en sammanfattning av de tester som har gjorts av gruppen. Eftersom testplanen inte var klar så har rapporten skrivits med hjälp av diskussioner under möten innan kundmötet och anteckningar från kundmötet.
\section{Enhetstest}
\subsection{Testaktivitet}
\subsection{Testresultat}
\subsection{Avvikelser}
\subsection{Utvärdering}
\section{Integrationstest}
\subsection{Testaktivitet}
\subsection{Testresultat}
\subsection{Avvikelser}
\subsection{Utvärdering}
\section{Systemtest}
\subsection{Testaktivitet}
\subsection{Testresultat}
\subsection{Avvikelser}
\subsection{Utvärdering}
\section{Acceptanstest}
\subsection{Pappersprototyp}
Tre stycken pappersprototyper hade tagits fram av gruppen och testades på kunden för att få feedback för att skapa en bättre prototyp. Två av dessa hade redan tidigare utvärderats och förbättrats i gruppen. Den tredje prototypen hade inte tidigare diskuterats i gruppen då denna skapades efter det diskussionstillfälle som var avsatt. Den var inte helt fullständig utan var främst tänkt att visa vissa alternativa lösningar på enskilda
\subsection{Testaktivitet}
Testet utfördes genom att Daniel från Region Östergötland fick i uppdrag att boka in en operation i vardera pappersprototyp. När han gjorde ett val så uppdaterade Martin gränssittet för att reflektera förändringen till följd av valet. Under testet så talade Daniel om vad han såg och hur han tänkte att han borde göra.
\subsection{Testresultat}
Utifrån de två prototyperna så listas de viktigaste punkterna som saknades i prototyperna nedan:
\begin{itemize}
\item Det kan vara bra med en kom ihåg mig ruta
\item Datum är bara siffror i rutor och därmed inte helt klart vad de innebär.
\item Det är inte helt klart att det är beslut som listas i listan med beslut.
\item Det finns en återställ knapp, vad gör den?
\item Vid bokning står det akut/elektiv vilket är fel.
\item Alla operationer kan inte utföras i alla salar.
\item Vissa parametrar kan ta mer plats än vad som givits.

\end{itemize}


\end{document}
